% ===================================================================
%  TABELLA N-RO 2 — Le corpore human. — Le recreation. — Le jocos.
%  Traduction 2026 d'apres texto/io, controlee sur texto/fr et
%  texto/es. Consigne : voir texto/ia/10-tabella-01.tex.
% ===================================================================

%%K t02-apar-1 apar
\VUsaut{4.0mm}
\VUtitre{15.0pt}{60}{TABELLA N\textsuperscript{o} 2}
\VUsaut{2.0mm}
\VUtitre{11.4pt}{0}{Le corpore human. --- Le recreation.}
\VUtitre{11.4pt}{0}{Le jocos.}

%%K t02-tit-0 sub
\VUtitre{13.2pt}{0}{Le corpore human.}
\VUsaut{2.4mm}
\VUcentre{10.2pt}{0}{\textit{(Simple lection de scientia natural.)}}
\VUsaut{3.0mm}

%%K t02-tit-1 sub pos=t02-01-1
\VUpk{10.2pt}{40}{Le capite.}
\VUsaut{3.0mm}

%%K t02-01-1 p
1. --- Le principal partes del corpore human es : le \VUgras{capite}, le
\VUgras{trunco} e le \VUgras{membros}.

%%K t02-02-1 p
2. --- Le capite comprende duo partes : le \VUgras{cranio}, que contine le
\VUgras{cerebro} e que coperi le \VUgras{capillatura}, e le \VUgras{facie}.

%%K t02-03-1 p
3. --- Sur nostre designo nos vide ni le cranio ni le cerebro; ma nos vide
multo clarmente le importante partes del facie.

%%K t02-04-1 p
4. --- Ecce primo le \VUgras{fronte}, que indica un potente intelligentia, un
pensata sempre agente. Le \VUgras{tempores} es multo libere. Le \VUgras{oculo}
es vivace e large aperte. Un dense \VUgras{supercilio} e longe \VUgras{cilios}
lo protege.

%%K t02-05-1 p
5. --- Ecce ancora le \VUgras{naso} e le \VUgras{narines}. Le naso nos servi
pro flairar e pro respirar. A cata latere del naso il ha le \VUgras{genas}, plus
lontano le \VUgras{aures}. Le parte inferior del facie consiste in le
\VUgras{bucca} e le \VUgras{mento}.

%%K t02-06-1 p
6. --- Nos non los vide multo ben, proque le \VUgras{mustachio} e le
\VUgras{barba lateral} nos los cela. Ma nos sape, que le bucca ha le duo
\VUgras{labios}, le \VUgras{maxilla superior} e le \VUgras{maxilla inferior} con
lor \VUgras{dentes}, e le \VUgras{lingua}. Per le bucca nos parla e mangia. Per
le lingua nos gusta le alimentos.

%%K t02-07-1 p
7. --- Le oculo, le aure, le naso e le bucca es, como le digitos, le organos del
cinque sensos : le \VUgras{vista}, le \VUgras{audito}, le \VUgras{olfacto}, le
\VUgras{gusto}, le \VUgras{tacto}.

%%K t02-08-1 p
8. --- Le capite es dunque un del plus interessante partes del corpore human.

%%K t02-tit-2 sub
\VUsaut{4.4mm}
\VUpk{10.2pt}{40}{Le trunco.}
\VUsaut{3.0mm}

%%K t02-09-1 p
9. --- Le \VUgras{collo} junge le capite al trunco. Illo comprende le
\VUgras{gorga} e le \VUgras{nucha}.

%%K t02-10-1 p
10. --- Le trunco comprende anque multe organos essential. In le \VUgras{pectore}
il ha le \VUgras{pulmones}, per le quales nos respira, e le \VUgras{corde} que
es le centro del vita. Quando le corde cessa batter, le esser vivente mori.

%%K t02-11-1 p
11. --- Le \VUgras{costas} mantene le pectore.

%%K t02-12-1 p
12. --- Le altere partes del trunco es le \VUgras{flanco}, le \VUgras{dorso}, le
\VUgras{spatulas} e le \VUgras{abdomine} (ventre). Nos nos corica sur le flanco
o sur le dorso pro dormir, nos porta le cargas sur nostre spatula.

%%K t02-13-1 p
13. --- In le abdomine il ha le \VUgras{stomacho} e le \VUgras{intestinos}.
Illos nos servi pro digerer le alimentos.

%%K t02-tit-3 sub
\VUsaut{4.4mm}
\VUpk{10.2pt}{40}{Le membros.}
\VUsaut{3.0mm}

%%K t02-14-1 p
14. --- Nos ha quatro \VUgras{membros} : duo \VUgras{bracios} e duo
\VUgras{gambas}. Per le bracios nos prende, tene, leva, collige le objectos; per
le gambas nos sta, marcha e curre.

%%K t02-15-1 p
15. --- Le \VUgras{spatula} junge le bracio al corpore. Un multo robuste
\VUgras{musculo}, le biceps, move le bracio que le \VUgras{cubito} junge al
\VUgras{antebracio}. Le \VUgras{carpo} constitue le articulation del
\VUgras{mano}, que fini per le \VUgras{digitos} e le \VUgras{ungues}. Sur le
mano nos distingue un \VUgras{vena} (e un arteria) in le qual flue le
\VUgras{sanguine}.

%%K t02-16-1 p
16. --- Le partes del gamba es : le \VUgras{femore}, le \VUgras{geniculo} e le
\VUgras{poplite}, le \VUgras{sura}, del qual le \VUgras{carne} es coperite per
le \VUgras{pelle}, le \VUgras{malleolo} e le \VUgras{pede}. Le \VUgras{ossos}
del gamba es multo grosse e multo forte. Al extremitate del pede il ha le
\VUgras{digitos del pede}. Le altere partes es le \VUgras{planta} e le
\VUgras{talon}.

%%K t02-17-1 p
17. --- Tal es le quasi complete description del corpore human.

%%K t02-17-2 p
\textit{Nota.} --- \textit{Con le adjuta del expression : « io ha dolor a... (le
capite etc.) » le professor potera reusar tote iste vocabulario del puncto de
vista del bon o mal} \VUgras{stato corporal.}

%%K t02-tit-4 sub
\VUsaut{5.0mm}
\VUpk{10.2pt}{40}{Gymnastica.}
\VUsaut{2.4mm}
\VUcentre{10.2pt}{0}{\textit{(Narration de un del alumnos.)}}
\VUsaut{3.0mm}

%%K t02-18-1 p
18. --- Pro esser flexibile, agile e san, nos debe exercer nostre gambas e
nostre bracios. Pro isto nos juga e practica le \VUgras{gymnastica}.

%%K t02-19-1 p
19. --- Durante le septimana, iste duo exercitios ha loco in le corte del lyceo
(vide le tabella n-ro 1). Ma le jovedi e le dominica nos va sur un grande
\VUgras{prato}, ubi nos ha spatio pro currer e amusar nos.

%%K t02-20-1 p
20. --- Nostre maestros e nostre parentes alcun vices nos accompania illac; ma
un \VUgras{professor de gymnastica} \textsuperscript{(12)} ja dirige le
exercitios.

%%K t02-21-1 p
21. --- Sur le prato, al sinistra del entrata, se trova le \VUgras{apparatos
gymnastic} \textsuperscript{(1)}, nos ascende per le \VUgras{scala}
\textsuperscript{(2)}, nos nos suspende capite in basso al \VUgras{anellos}
\textsuperscript{(3)} e al \VUgras{trapezio} \textsuperscript{(4)}, nos nos
hissa per le manos usque al summitate del \VUgras{scala de corda}
\textsuperscript{(5)} o del \VUgras{corda a nodos} \textsuperscript{(6)}.

%%K t02-22-1 p
22. --- Intertanto, nostre camerados salta super le \VUgras{cavallo de ligno}
\textsuperscript{(7)} o le \VUgras{barras parallel} \textsuperscript{(11)}.
Alteres \VUgras{boxa} \textsuperscript{(8)} o \VUgras{scherma}
\textsuperscript{(9)} con masca e per un \VUgras{floretto}. Finalmente,
alteres \VUgras{leva halteres} \textsuperscript{(10)}.

%%K t02-tit-5 sub
\VUsaut{4.6mm}
\VUpk{10.2pt}{40}{Le jocos.}
\VUsaut{3.0mm}

%%K t02-23-1 p
23. --- Circa le gymnastas, pueros e pueras juga a diverse \VUgras{jocos}. Le
pueras se amusa con lor \VUgras{cerculo} \textsuperscript{(14)}; illas salta al
\VUgras{corda} \textsuperscript{(13)} o invita lor fratres e cosinos a un
partita de \VUgras{croquet} \textsuperscript{(15)}. Illas se delecta a repeller
le \VUgras{ballon} del adversario con lor \VUgras{martelletto de ligno}, al
momento quando ille pensa passar facilemente le arco!

%%K t02-24-1 p
24. --- Le pueros prefere le jocos plus muscular. Con placer illes juga al
\VUgras{quilias} \textsuperscript{(16)} que illes abatte habilemente per un
\VUgras{bolla} \textsuperscript{(17)}. Illes anque non dedigna jugar al
\VUgras{trompo} \textsuperscript{(18)} e al \VUgras{bilboquet}
\textsuperscript{(19)} o mesmo al \VUgras{rana} \textsuperscript{(20)}. Ma lor
jocos preferite es ben le \VUgras{billas} \textsuperscript{(21)} e le
\VUgras{pilota contra le muro} \textsuperscript{(22)}, proque le muro del
\VUgras{serra} \textsuperscript{(26)} es multo apte pro facer resaltar le legier
pilota.

%%K t02-25-1 p
25. --- Le parve Johannes, qui va sur su \VUgras{trampolos}
\textsuperscript{(24)}, es multo habile e sape multo ben conservar le
equilibrio. Presso ille, alcun camerados juga al \VUgras{celar-se}
\textsuperscript{(25)} in ille serra e detra le \VUgras{haga}. Alteres prefere
le \VUgras{colpo a divinar} \textsuperscript{(28)} o le \VUgras{volante}
\textsuperscript{(29)} que vola de un \VUgras{racchetta} al altere.

%%K t02-noto-1 noto
\VUnotes{143.2mm}{%
(*) Le jocos de infantes o de pueros ha sovente nomines purmente national. Nos
ha agite pro nominar los in Ido le melio possibile, sia inspirante nos del
action mesme, per le qual illos es characterisate, sia eligente le nomine
national in iste o ille paises, quando illo non es troppo injuste. Ex. : le
\VUgras{joco del barril} (in germano e in francese) es le \VUgras{joco del rana}
\textsuperscript{(20)} (in anglese) in le qual le jugatores se effortia jectar
lor discos rotunde tra le bucca del rana de metallo que sta con le bucca aperte
sur le cassa (barril) traforate. Le \VUgras{salto super le spatulas}
\textsuperscript{(30)} differe del \VUgras{salto super le dorso} solmente per
isto : pro saltar super le capite, le jugatores appoia le manos sur le spatulas
de lor partenario, durante que in le « \VUgras{salto super le dorso} » illes
appoia le manos solmente sur su dorso. --- O anque alcunes del participantes es
inclinate durante que le alteres saltara super lor dorso appoiante le manos sur
le spatula; le primes debe tolerar le colpo sin flecter se, le secundes debe
saltar sin perder le equilibrio (Vide le tabella mesme).%
}

%%K t02-26-1 p
26. --- In le medio del prato il ha duo arbores. Al pede de un del duo, septe o
octo pueros ha organisate un partita de \VUgras{salto super le spatulas}
\textsuperscript{(30)} (*). Non in tote le paises on cognosce iste joco; ma
totes sape lo que es le \VUgras{salto super le dorso}, al qual le pueros non
juga iste vice.

%%K t02-tit-6 sub
\VUsaut{5.0mm}
\VUpk{10.2pt}{40}{Le jocos in aere libere.}
\VUsaut{3.4mm}

%%K t02-27-1 p
27. Le \VUgras{joco del cieco} \textsuperscript{(31)} es anque multo amusante;
pueras e pueros alternatemente pone le \VUgras{banda sur le oculos} e sasi le
maladroites qui se lassa capturar.

%%K t02-28-1 p
28. --- In le medio del prato, ecce un joco multo francese ma un pauco violente
: illo es le \VUgras{ursa} \textsuperscript{(32)} multo plus bruente que le tanto
tranquille joco del \VUgras{cerf-volante} \textsuperscript{(33)} o mesmo que le
joco anglese del \VUgras{tennis} \textsuperscript{(34)}.

%%K t02-29-1 p
29. --- Iste ultime sport es le gaudio del grande alumnos. Nostre professores
mesme lo practica con placer. Illo require un grande habilitate con agilitate e
illo me place multo. Io lassa al pueras le jocos del \VUgras{balancia}
\textsuperscript{(35)}.

%%K t02-30-1 p
30. --- Io non ama un altere joco anglese que forsan devenirea periculose.

%%K t02-30-1b p
Io intende parlar del \VUgras{football} \textsuperscript{(36)} que allegra mi
fratre plus vetule. Io prefere nostre vetule \VUgras{joco del barras}
\textsuperscript{(38)} tanto ben hygienic e vermente minus brutal.

%%K t02-tit-7 sub
\VUsaut{4.6mm}
\VUpk{10.2pt}{40}{Cyclismo e joco de pilota.}
\VUsaut{3.0mm}

%%K t02-31-1 p
31. --- Anque con placer io circumvade le prato per \VUgras{bicycletta}
\textsuperscript{(37)}.

%%K t02-32-1 p
32. --- Durante le anno proxime io apprendera le \VUgras{equitation}
\textsuperscript{(39)}. Tanto belle es un cavallo qui gratiosemente passa o
trotta, e qui al galopo salta super le obstaculos!

%%K t02-33-1 p
33. --- In estate, nos apprende le \VUgras{arte de natar}
\textsuperscript{(40)}. Nos con delecto plongia e nos bania in le clar e fresc
aqua del riviera. Alcun vices finalmente nos lancea un \VUgras{ballon}
\textsuperscript{(41)} de papiro que majestosemente ascende in le atmosphera,
tosto que illo es plenate de aere calide.

%%K t02-34-1 p
34. --- Il ha anque multe altere jocos, ma io non finirea enumerar los e on se
rende conto secundo iste resumito, que le jovedis, on non se enoia multo sur
nostre belle prato.

%%K t02-34-2 p
N.-B. --- \textit{Le professor completara facilemente iste capitulo; describente
plus detaliatemente cata un del plus importante jocos.}
\VUsaut{6.0mm}
\VUfilet{22mm}
